\chapter{Experimental Setup}
\label{ch:setup}

\section{Acoustic frontend and duration policy}
\label{sec:setup-frontend}

\subsection{AT-ADD Frontend and Duration Normalization}
\label{sec:duration-norm}

The baseline frontend follows the requirements of AT-ADD Track 2, which evaluates type-agnostic detection across speech, environmental sound, singing and music~\cite{xie2026atadd}. Audio is converted to mono, resampled to 24\,kHz, and represented using a fixed four-second input window. Longer recordings are randomly cropped during training and centre-cropped during evaluation, while shorter recordings are extended using repeat-padding.

This duration policy was motivated by the heterogeneous duration distribution of AT-ADD. The generator-level statistics used in this study show that many music and sound sources have approximately 10-second recordings, whereas numerous speech generators have median durations close to or below four seconds; for example, the median duration is 3.66\,s for ChatTTS and 3.51\,s for DiffGANTTS. Thus, a substantial portion of speech data requires extension to reach the fixed input duration. Zero-padding would introduce a deterministic trailing-silence region that is unrelated to the acoustic content and could provide a potential duration cue to the classifier. Repeat-padding avoids this particular silence pattern, although the repetition boundary may itself introduce an acoustic discontinuity.

The same frontend is used for E1, E2 and E3 so that the baseline comparison changes the detector architecture or learned representation without simultaneously changing the input preprocessing.

\section{Compute and software}
\label{sec:compute}

Experiments were implemented in PyTorch using CUDA-accelerated training with automatic mixed precision. Audio was decoded using SoundFile and resampled with librosa; MP3 coding for degradation simulations uses FFmpeg. Random seeds were fixed (seeds 42, 1337, and 2024 for the three-seed baseline experiments; seed 42 for exploratory downstream models).

\section{Training budget}
\label{sec:budget}

The clean baseline candidate models (E1, E2, E3) were trained on the 60{,}000-sample training manifest with a standardized maximum budget of 25 epochs. In all cases the model checkpoint was selected based on the highest Macro-F1 on the 22{,}536-sample development manifest. Although the training run for E1 seed~1 was permitted to run for 35 epochs, its best validation checkpoint occurred early at epoch~7 (before epoch~10); its reported metrics therefore faithfully reflect the standardized 25-epoch budget. Downstream exploratory models (FT-RCL and STAR-Edge) were trained under a 20-epoch budget on the MobileNetV3-Small backbone with early stopping (patience 5).

\section{Hyper-parameters}
\label{sec:hyper}

Table~\ref{tab:hyper} details the training configuration for the baseline selection experiments, distinguishing shared signal-processing settings from model-specific optimiser schedules. All three baseline candidates incorporate gradient clipping ($\text{MAX\_GRAD\_NORM}=5.0$) where gradients are unscaled from mixed precision prior to clipping.

\begin{table}[htbp]
\centering
\caption{Training configuration for clean baseline candidate experiments and downstream models.}
\label{tab:hyper}
\small
\begin{adjustbox}{max width=\textwidth}
\begin{tabular}{@{}ll@{}}
\toprule
Parameter & Value \\
\midrule
\multicolumn{2}{@{}l}{\textit{Common baseline settings (E1, E2, E3)}} \\
Input audio & 4\,s, 24\,kHz, 80-bin log-mel (window 600, hop 240, FFT 1024) \\
Normalisation & Training-set global z-score ($\mu \approx -9.128$, $\sigma \approx 6.774$) \\
Optimiser & AdamW, weight decay $10^{-4}$ \\
Batch size & 32 \\
Epochs budget & 25 epochs \\
Checkpoint criterion & Maximum development Macro-F1 \\
Gradient clipping & Norm clipping with $\text{MAX\_GRAD\_NORM} = 5.0$ (unscaled prior to clipping) \\
Random seeds & 3 independent seeds (42, 1337, 2024) \\
Loss function & Class-weighted cross-entropy ($\omega_0 = 1.875$, $\omega_1 \approx 0.682$) \\
\midrule
\multicolumn{2}{@{}l}{\textit{Model-specific settings}} \\
E1 (MobileNetV3-Small) & Scratch, 1{,}519{,}618 params; uniform lr $10^{-3}$; ReduceLROnPlateau (factor 0.5, patience 3) \\
E2 (MobileNetV3-Small + GRU) & Scratch, 1{,}198{,}082 params; uniform lr $10^{-3}$; ReduceLROnPlateau (factor 0.5, patience 3) \\
E3 / M0 (PANNs MobileNetV2) & AudioSet-pretrained, 3{,}537{,}250 params; frozen pretrained BN stats; \\
 & Layer-wise lr: features $10^{-5}$, FC1 $5\times 10^{-5}$, fresh BN0 $3\times 10^{-4}$, classifier $3\times 10^{-4}$; \\
 & ReduceLROnPlateau (factor 0.5, patience 3, min lr $10^{-6}$) \\
\midrule
\multicolumn{2}{@{}l}{\textit{Downstream models (exploratory MobileNetV3-Small runs)}} \\
FT-RCL (M3) & $\beta=0.5$, $\gamma_{\mathrm{pc}}=1.0$, $\gamma_{\mathrm{ftsc}}=0.5$, $K=8$; type weights $\alpha_t$ (speech/singing 1.0, sound 0.5, music 0.3); \\
 & 20 epochs, seed 42, $J$-score checkpoint selection \\
STAR-Edge (S) & ABR ($m=0.05$), STC, CAD; 20 epochs, seed 42, $J$-score checkpoint selection \\
\bottomrule
\end{tabular}
\end{adjustbox}
\end{table}

\section{Catalogue of experiments}
\label{sec:cat}

Every experiment evaluated in this study is listed in Table~\ref{tab:catalogue} with a unified identifier.

\begin{table}[htbp]
\centering
\caption{Experiments reported in this work.}
\label{tab:catalogue}
\small
\begin{adjustbox}{max width=\textwidth}
\begin{tabular}{@{}p{1.1cm}p{6.3cm}p{6.6cm}@{}}
\toprule
ID & Experiment & Purpose and Setting \\
\midrule
P4 & Re-implementation of published MobileNetV3 system~\cite{pham2024deepfake} & ASVspoof~2019 LA architectural reference \\
\midrule
E1 & MobileNetV3-Small scratch baseline & Compact CNN candidate, 3 seeds (42, 1337, 2024), 25 epochs \\
E2 & MobileNetV3-Small + 1-layer GRU & Temporal sequence candidate, 3 seeds, 25 epochs \\
E3 & PANNs MobileNetV2, AudioSet-pretrained & Pretrained audio representation candidate, 3 seeds, 25 epochs \\
M0 & Selected clean baseline (E3 configuration) & Selected clean reference model for downstream benchmarks \\
\midrule
M3 & FT-RCL (full objective) & Dual-view coherence on MobileNetV3-Small backbone, seed 42 \\
S & STAR-Edge & Extension on MobileNetV3-Small backbone, seed 42 \\
R & Degraded-audio pilot on MobileNetV3-Small and FT-RCL & Fake music clips only (not valid) \\
C & Shortcut-cue analysis & About 6{,}000 training and 6{,}000 evaluation clips \\
D & Duration and padding statistics & All 126{,}825 clips \\
\bottomrule
\end{tabular}
\end{adjustbox}
\end{table}

Modes M1 and M2 of Table~\ref{tab:modes} are defined but were not run.
